\clearpage
\fancyhead[L]{\small IRIS.}
\SourceLeaf{480}
\begin{center}
\rule{0.55\linewidth}{0.4pt}\\[2pt]
Astronomische Nachrichten. Band XXVI. Nr. 616. S. 241 u. 255. 8. November 1847.\\[2pt]
\rule{0.55\linewidth}{0.4pt}
\end{center}
\begin{center}[\textsc{Gauss} an \textsc{Schumacher}.]\quad Göttingen, 11. Oktober 1847.\end{center}

Nach einer Unterbrechung von vollen drei Wochen kann ich Ihnen heute doch noch eine Meridianbeobachtung der \emph{Iris} mitteilen. Der Planet war allerdings wegen seiner geringen Helligkeit sehr schwer zu beobachten.

\begin{center}
\begin{tabular}{ll}
\multicolumn{2}{l}{1847 Oct. 10.\quad $6^{\rm u}\ 42^{\rm m}\ 2^{\rm s}6$ M. Z.} \\
Rektasz. & $299^\circ\ 13'\ 32''7$ \\
Südl. Dekl. & $14^\circ\ 19'\ 38''6$ \\
\end{tabular}
\end{center}

\textsc{Goldschmidts} Elemente stimmen noch sehr gut. Von der \emph{Sonnenfinsternis} konnte ich nur das Ende beobachten.

\begin{center}
1847 Oct. 8.\quad $21^{\rm u}\ 39^{\rm m}\ 33^{\rm s}7$ Mittl. Zeit.\\
Professor \textsc{Goldschmidt} $2^{\rm s}4$ später.
\end{center}

Außerdem habe ich den Eintritt eines Flecks aufgezeichnet, also den Beginn des Eintritts oder die erste Berührung, um $19^{\rm u}\ 3^{\rm m}\ 40^{\rm s}8$. Die Luft flimmerte dabei so stark, dass das deutliche Sehen sehr beeinträchtigt wurde. Professor \textsc{Goldschmidt} hatte $3^{\rm s}7$ weniger.

In einem späteren Schreiben fügte der Geheime Hofrat noch folgende Meridianbeobachtungen der \emph{Iris} hinzu:

\begin{center}
\setlength{\tabcolsep}{4pt}
\begin{tabular}{l|r|r|r}
1847 & Mittlere Zeit & Rektaszension & Südl. Dekl. \\\hline
Oktober 11 & $6^{\rm u}\ 39^{\rm m}\ 0^{\rm s}6$ & $299^\circ\ 27'\ 3''6$ & $14^\circ\ 18'\ 39''3$ \\
\hfill 13 & $6\ 33\ 1.5$ & $299\ 55\ 18.2$ & $14\ 16\ 11.0$ \\
\hfill 14 & $6\ 30\ 4.2$ & $300\ 9\ 59.4$ & $14\ 14\ 44.5$ \\
\end{tabular}
\end{center}

\clearpage
\hypertarget{native-480-october21}{}
\begin{center}[\textsc{Gauss} an \textsc{Schumacher}.]\quad Göttingen, 21. Oktober 1847.\end{center}

Ich kann Ihnen noch drei Meridianbeobachtungen der \emph{Iris} schicken. Ich führe hier aber auch die Beobachtungen vom 13. und 14. noch einmal auf, weil sich die Deklinationen infolge der neuen Nadirbestimmungen geringfügig geändert haben und ich nicht weiß, ob die Änderung bei den Zahlen, die ich Ihnen früher mitgeteilt habe, schon berücksichtigt war. Die Angaben vom 10. und 11. sind davon nicht betroffen.

\begin{center}
\setlength{\tabcolsep}{4pt}
\begin{tabular}{l|r|r|r}
1847 & Mittlere Zeit & Rektaszension & Südl. Dekl. \\\hline
Oktober 13 & $6^{\rm u}\ 33^{\rm m}\ 1^{\rm s}5$ & $299^\circ\ 55'\ 18''1$ & $14^\circ\ 16'\ 11''4$ \\
\hfill 14 & $6\ 30\ 4.2$ & $300\ 9\ 59.4$ & $14\ 14\ 45.1$ \\
\hfill 15 & $6\ 27\ 8.2$ & $300\ 25\ 1.5$ & $14\ 13\ 17.2$ \\
\hfill 16 & $6\ 24\ 13.9$ & $300\ 40\ 27.1$ & $14\ 11\ 50.3$ \\
\hfill 17 & $6\ 21\ 20.8$ & $300\ 56\ 12.1$ & $14\ 10\ 9.1$ \\
\end{tabular}
\end{center}

Seitdem war es um die Kulminationszeit stets trüb. Das habe ich umso mehr bedauert, als es mich daran hinderte, den Planeten in sehr enger Konjunktion mit einem Stern 9. Größe aus \textsc{Bessels} Zonen zu sehen (20, 101 \textsc{Weisse}, dem der Planet am 18. bis auf fast eine halbe Minute nahe gekommen sein muss). Übrigens war der Planet wegen seiner sehr geringen Helligkeit stets sehr schwer zu beobachten. Ich kann aber nicht sagen, dass die große Nähe des Mondes am 17. die Beobachtung schwieriger gemacht hat als an den früheren Abenden.